\setcounter{section}{1}
\section{Preliminaries}
\setcounter{theorem}{15}
%

In this section we establish notation and we formally state basic facts and definitions which
we will need in later sections.



\begin{definition}[3-dimensional matching~\cite{K72}]\label{def:3D-matching}
	The \emph{3-dimensional matching} problem is defined as follows:
	\begin{itemize}
		\item[Input:] a set $U \subseteq [n] \times [n] \times [n]$, representing the edges of
		a tripartite, 3-uniform hypergraph.
		\item[Output:] YES, if there is a set of hyperedges $W \subseteq U$ such that $|W| = n$ and
		no two elements of $W$ agree in any coordinate (that is, they form a matching in this
		hypergraph). NO, if there is no such set.
	\end{itemize}
\end{definition}

\begin{theorem}[NP-completeness of 3-dimensional matching~\cite{K72}]\label{thm:3D-matching-NP-complete}
	The 3-dimensional matching problem is NP-complete.
\end{theorem}

\subsection{Basic facts from algebraic complexity}

We now give basic facts that from algebraic complexity which we will use in the next sections.

The next proposition shows that homogeneous components of low degree
of an arithmetic circuit can be efficiently computed, with a small blow-up
in circuit size and without the use of any extra constants. This proposition was
originally proved by Strassen in~\cite{S73} and its proof can be found
in~\cite[Theorem 2.2]{SY10}.
In the following proposition, given a polynomial $p(\bx)$, we denote its degree-$d$ homogeneous component by $H_d[p(\bx)]$.

\begin{proposition}[Efficient computation of homogeneous components]\label{prop:hom-components}
	Given a circuit $\cC(\bx)$ of size $s$, then for every $r \in \N$
	there is a homogeneous circuit $\Psi(\bx)$ of size $O(r^2 s)$
	computing $H_0[\cC(\bx)], H_1[\cC(\bx)], \ldots, H_r[\cC(\bx)]$.
	Moreover, the constants used in the computation of the components $H_i[\cC(\bx)]$
	are a subset of the coefficients used in the computation of $\cC(\bx)$.
%
%
\end{proposition}

The next theorem, proved by~\cite[Theorem 4.10]{AB03} gives us a randomized polynomial time
algorithm to test whether an algebraic circuit of polynomial size, with rational
coefficients, is identically zero. Another randomized algorithm easily follows
from~\cite[Lemma 2]{S79}, when adapted for polynomials with rational coefficients.

\begin{theorem}[PIT for poly-sized circuits~\cite{AB03}]\label{thm:PIT-BPP}
	Let $P(\bx) \in \Q[\bx]$ be a polynomial in the variables $\bx = (x_1, \ldots, x_n)$, with
	each variable $x_i$ having degree bounded by $d_i$, and whose coefficients are rationals with
	bit complexity bounded by $B$. If $P(\bx)$ is given as an arithmetic circuit of size $s$, then
	there exists a randomized algorithm running in time ${\rm poly}(n, s, \log(B), 1/\epsilon)$
	and using $O\left(\sum_{i=1}^n \log(d_i) + \log(B) \right)$ random bits which tests whether $P(\bx)$ is identically zero.
	If $P(\bx)$ is the identically zero polynomial then the algorithm always succeeds.
	Otherwise, it errs with probability at most $\epsilon$.
\end{theorem}
